\chapter{}

1911 年 10 月 10 日，武昌楚望台的枪声一响，各地纷纷响应，遂成燎原之势。统治中国 268 年的清政府在革命呼声中迅速坍塌，“中华民国”就这样在一片仓促的欢呼声中诞生。已经领导了十次起义的孙中山此时正在美国北部哥罗拉多州筹款，听到这个消息是几天以后，当时他也十分惊讶，胜利得来竟在不经意间。他经过再三考虑，认为自己当前的主要工作，不在“疆场之上”，而在“樽俎之间”，他希望通过外交活动，断绝清政府的后援。他没有立即回国，而是去了伦敦。在他的心目中，得到列强承认，筹到足够的金钱，才是当务之急。这一着棋孙中山显然没有走好，当然那时大家都不清楚新生的革命政权最需要什么。

资产阶级革命党人是在民族危机的严重关头，思想准备十分不足的情况下走上革命道路的。自从 20 世纪初开始，他们的革命理由，主要是民族危亡和清廷统治。他们认为亡国灭种危机是清朝的反动卖国政府造成的，清廷放手卖国的原因就在于它是一个“异族”的朝廷，所以对汉族祖先家财毫不吝惜。因此对于怎样革除封建统治，建立什么样的国家，思考不足。他们提出“驱除鞑虏”的口号，实质上也包含着对外反对帝国主义的侵略、挽救民族危亡，对内反对封建压迫、种族歧视的内容，因此它能够迅速为广大民众所接受。但对于“平均地权，建立合众国”，实际上想的不多。胡汉民后来在总结经验时说：“同盟会未尝深植其基础于民众，民众所接受者，仅三民主义中之狭义的民族主义耳。正惟‘排满’二字之口号，亟简明切要，易于普遍全国，而弱点亦在于此。民众以为清室退位，即天下事大定，所谓‘民国共和’则仅取得从来未有之名词而已。至其实质如何都非所向。”武昌起义后，不仅很多一般的革命党人，以为只要清帝退位，共和政府成立，汉人做了大总统，就算是革命成功了，就连孙中山、黄兴这样的革命领袖也是这样认为的。正因为如此，1911 年 12 月 25 日，孙中山从国外归来，1912 年1 月1 日孙中山在南京就任临时大总统才 3 个月，孙中山就把政权交给了袁世凯。资产阶级革命派的领袖、共和国的创始人孙中山让位于清廷旧臣，这是关系辛亥革命成败的重大问题，就这样了结了。当然这不是某一个人的主观意愿，其背后复杂深刻的社会历史背景，才是我们需要反思的。

首先必须注意到，武昌起义后才几天，1911 年10 月21 日的《民立报》以《欧洲关于中国革命之电报》为题的新闻综述中说：“《每日镜》、《伦敦晚报》及其他各报宣言孙逸仙已选袁世凯为第一总统。此间舆论极赞成袁世凯联合革命党，并望孙勿念旧日之恨，袁当有以助其成功，云云。”紧跟着便出现了无人署名的《鄂人致袁世凯书》，和革命党人《敬告袁项城》，这些都不是无穴来风。当时的长江流域是英国的势力范围，英国在中国的贸易总额包括香港在内，超过其他各资本主义国家在华贸易的总和，所以英国最需要一个符合他们在华利益的总统。孙中山在国外时，就已经听到一种舆论，即如果争取到袁世凯拥护共和制度，可以让袁出任民国总统。孙中山原来对袁世凯的印象并不怎么好，觉得此人反复无常“狡猾善变”，不太靠得住。但他又希望避免流血，尽早实现革命目标，只要推翻清政府，废除帝制，即使是袁世凯出来当总统，也未尝不可。1911 年 12 月 25 日，孙中山从国外归来，面对着第一次各省都督代表会议通过的“若袁世凯反正，当公举为临时大总统”这样的决议，他也承认这个事实。孙中山当选为临时大总统后，主张“让位”的空气仍然笼罩着革命党人，包括孙中山身边的一些重要人物，如黄兴、汪精卫、胡汉民等人，都表示赞成让位。汪精卫甚至讽刺孙中山说：“你不赞成议和，难道是舍不得总统的职位吗？”革命党人的二号人物、担任临时政府陆军总长的黄兴，对袁世凯既有顾虑，又存幻想。黄兴说，袁世凯是一个奸猾狡诈、胆大妄为的人，如能满足他的欲望，他可以帮助我们推翻清朝。否则，他也可以像曾国藩替清朝出力搞垮太平天国一样来搞垮革命。只要他肯推翻清朝，我们给他一个民选的总统，任期不过几年，可以使战争早停，民众早过太平日子，岂不很好？黄兴的这种看法，在当时革命党人中是很有代表性的，也完全符合当时孙中山的思想状态。正是由于革命党人过分强调满汉对立，简单地宣传“排满”，这就使他们不仅没有把汉族的官僚和军阀当作革命对象，反而把他们当作可以争取的同胞兄弟。所谓“论地位则为仇雠”，“论情谊则为兄弟”。只要他们站到“反满”的行列中来，即可“离仇雠之地位而复为兄弟”。这种长期而反复的宣传，在革命党人和一般的民众中自然产生了广泛的影响。武昌起义后，许多人继续强调满汉矛盾，接受甚至拥戴清朝的督抚宣布独立，正是这种思想指导下的结果。“举袁”方针的提出，自然与这种指导思想是分不开的。袁世凯既为汉人，只要他站到“反满”的行列中来，他就可以“复为兄弟”。从这一点来讲，当时提出只要袁世凯反正即可举为大总统，不仅不足为怪，而且也是合乎逻辑的。

另外一个重要原因是，害怕革命战争的延长引起帝国主义的干涉他们看到帝国主义列强对袁世凯的支持是为利益，认为“举袁可以杜外人之干涉”。本来革命党是以反对帝国主义的侵略，救亡图存为出发点的，可是由于民族资产阶级自身的软弱，又看不到抗拒帝国主义的力量，他们不但不敢提出反对帝国主义的口号，相反恐惧帝国主义的干涉。他们小心翼翼地力图避免革命损害了帝国主义在中国的现实利益。早在同盟会的《对外宣言》中，就明确宣布承认清政府与列强所签订的一切不平等条约，“偿款外债照旧承认”，“所有外人之既得权利，一体保护”。武昌起义后，各省军政府都严格执行了《对外宣言》中所规定的各项承诺。他们急切地希望得到列强的承认，即使有一个外国人来电祝贺革命成功，也被说成是“新共和国外交之成功”。他们害怕革命会因帝国主义的干涉而遭受到太平天国那样的失败。这种害怕列强干涉的心理，可以看到自 1840 年以来失败，对中国民主资产阶级及知识精英自信心的打击是何等深刻。孙中山在得到武昌起义的消息后，绕道欧洲进行外交活动，到达伦敦后的第三天，即 11 月14 日，英国外交大臣格雷就让人转告孙，英国政府对袁世凯将给予尊敬，“所有外国人以及反满的团体都可能给予袁世凯以总统职位—如果他能驱逐满清并赞成共和。”孙中山则是把对英国外交的成败，看成是“可以举足轻重为我成败存亡所系”的。英国外交大臣公开表示支持袁世凯做大总统，自然使他不能不加以郑重考虑。所以，当孙中山接到上海已有议会的组织，将举黎元洪或袁世凯做大总统时，即于11 月16 日致电民国军政府说：“今闻已有上海议会之组织，欣慰。总统自当推定黎君，闻黎有推袁之说，合宜亦善。总之，随宜推定，但求早固国基。”也就是说，在他回国以前已经接受这样的安排了。12 月20 日，他归国途中经香港与胡汉民谈话时所说的一段话，可以看作他接受这个安排的一个说明：“革命军骤起，有不可向迩之势，列强仓卒，无以为计，故只得守向来局外中立之惯例，不事干涉。然若我方形势顿挫，则此事正未可深恃。戈登、白齐文之于太平天国，此等手段正多，胡可不虑？谓袁世凯不可信诚然。但我因而利用之，使推翻二百六十余年贵族专制之满洲，则贤于用兵十万。纵其欲继满洲以为恶，而其基础已远不如，覆之自易，故今日可先成一圆满之段落。”孙中山的这段话表明：他之所以采用“举袁”的方针，就是因为害怕革命战争的延长引起帝国主义的干涉，导致革命像太平天国那样的失败。

另外，资产阶级革命一大特点就是找钱，孙中山一生都在筹款。武昌起义后，尚在美国的孙中山即把财政问题视为革命的成败关键。他从美国到伦敦和巴黎活动的目的之一，即为进行贷款。可是，均遭拒绝。但孙中山并未因此放弃从列强取得贷款的幻想，因为他相信法国东方汇理银行总裁向他所说的“一旦民军建立起一个为全国所接受，为列强所承认之正规政府时，他们对于在财政上帮助革命党，将不表反对”。取得外国的贷款，一直是孙中山所希望的。为此在筹组临时政府考虑财政总长人选时，孙中山也从有利于取得外国借款的角度出发，以陈锦涛“曾为清廷订币制，借款于国际，有信用”，决定选用陈锦涛。陈上任后，被授予募筹款项以应财政紧迫需要的重大任务，前去上海谋求外国的贷款。可是，一直没有进展。但孙中山还没放弃向外国贷款的幻想。财政问题中军饷是最为紧迫，有的军官扬言：“军队乏饷即溃，到那时只好自由行动，莫怪对不住地方。”黄兴为军饷问题，急得天天跳脚。各省军政府亦同样存在着军饷问题，有的军政府为此还向临时政府伸手。严重的财政危机，迫使孙中山不得不放弃原拟采取的“若新政府借外债，则一不失主权，二不抵押，三利息轻”的立场，考虑由私人企业出面，接受条件极为苛刻、有损民族利益的外国贷款。拟将轮船招商局为抵押和同意中日合办汉冶萍煤铁公司，换取日本垄断资本提供大笔贷款。结果遭到舆论的激烈反对，只好废除了原已达成的初步协议。临时政府的这种举动，虽出于万不得已，但其名声却不能不受到损害。孙中山在致章太炎的信中说：“此事(指汉冶萍借款事)，弟非不知利权有外溢之处，其不敢爱惜声名，冒不韪而为之者，犹之天寒解衣付质，疗饥为急。先生等盖未知南京军队之现状也。每日到陆军部取饷者数十起。”“无论和战如何，军人无术使之枵腹。前敌之士，犹时有哗溃之势。弟坐视克兄(指黄兴)之困，而环视各省，又无一钱供给。”“至于急不择荫之实情，无有隐饰，则祈达人之我谅。”孙中山的这封信，十分清楚地表明他当时在财政方面所处的困境。南京临时政府一开始就陷入严重的财政困难中， 因为帝国主义在武昌起义后截夺了中国海关全部税收。中国海关税收虽早为帝国主义指定为偿付外债和赔款，但税款的保管和支付，完全由清政府委任的海关道或海关监督全权负责。外籍税务司的权力，仅限于征收关税。武昌起义爆发后，帝国主义为了维护它们在中国的侵略利益和完全控制中国的海关，立即开始攫取中国海关税款的活动。10 月15 日，总务税司安格联说：“让税款跑到革命党的库里去是不行的”，他下令给江汉关税务司苏古敦“将税款设法付入汇丰银行我的账内”。后来，经朱尔典提议，外国公使团竟然决议“把全部海关岁入置于总税务司的控制下”，由列强在上海组织专门机构负责接受这些税款。对此，各地的革命党人虽都先后进行了抗争，但由于害怕帝国主义的武装干涉，最后都不了了之。这样，武昌起义后，凡爆发革命的各通商口岸的海关税收，无一例外都落入列强的口袋。

从武昌起义到孙中山“让位”，时间不到半年。但辛亥革命毕竟结束了两千多年的封建帝制，对中国进入现代政治具有转折性影响，其后又经过“二次革命”“三次革命”，各种力量陆续登上政治舞台。尽管对老百姓来说不过是把头上的辫子剪掉了，社会生活没有发生多少改变，但在意识形态上破除了“三纲五常”，客观上为新文化运动扫清了障碍。同时这段历史已经把中国要取得民族独立、民主革命胜利所需要的种种主客观条件都展示出来了，资产阶级革命迷信资本惧怕列强、不相信人民不依靠人民的特性也都暴露得很充分。孙中山是很愿意总结经验教训的，他后来终于认识到把政权拱手让给袁世凯是一个历史性的错误。他沉痛地写道：“我的辞职是一个巨大的政治错误，它的政治后果正像在俄国如果让高尔察克、尤登尼奇或弗兰格尔跑到莫斯科去代替列宁而就会发生的一样。”

然而到了 1980 年代以后，学术界却突然刮起一股否定辛亥革命鼓吹袁世凯的怪异之风。其代表性的论调是“革命党不如维新党，维新党不如保皇党，保皇党不如慈禧老佛爷”，“仓促的革命打断了晚清的温和的政改之路”，意思是慈禧也是支持改革的，如果坚持改革社会进步要大得多、代价要小得多。此论调的代表人物是历史学教授袁伟时，他不但否定辛亥革命，而且鼓吹“晚清政改”，“清末新政在几个领域的大改革，是非常深刻的变革。”甚至认为八国联军侵华是因为“违反国际法，由此遭到英法联军的正义惩罚，火烧圆明园乃咎由自取；义和团违反国际法，八国联军乃出堂堂之师，庚子之变实乃祸由己出！”这位老先生频频发出惊人之语，经由凤凰网腾讯网的接力传播，居然辛亥革命是干什么的都搞不清楚了，他说：“就辛亥革命而言，主要不是解决经济问题，是要解决政治上的专制制度，要求民主，保障公民的自由、平等。这是近代革命的核心内容，也是辛亥革命的基本诉求。”如此一来，“驱除鞑虏，恢复中华，平均地权，建立合众国”不算数了，孙中山还要重新培训才能上岗。

所谓听话听音锣鼓听声，他这样说已经跟历史学无关，而是要表达自己的现实政治诉求，那就是中国尽快实行宪政。他对宪政的向往已经到了迷恋：“以各省咨议局和商会为中心的民间社会相当强大，他们的独立性很强，几次请开国会运动就是他们领导的。”如此美化宪政的袁伟时不会不知道，民国北洋政府这一段，中国先后有大大小小300 多个政党，不是没有搞过宪政。大体上是袁世凯4 年、皖系 4年、直系4 年、奉系 4 年，袁世凯、段祺瑞、曹锟、吴佩孚、张作霖依次各唱了几年的宪政戏。这 16 年，“民犹是也，国犹是也，无分南北；总而言之，统而言之，不是东西。”连选总统都可以收钱投票的国会，3000 元一票5000 元一票，自然做不成代议机关。宪政丑态百出成为一个笑话，张作霖当时的一句名言是：“管你吃，管你穿，不听话怎么能行？汉高祖刘邦约法三章，我只一章，不听话就枪毙。”那么袁伟时为什么还要这样死乞白赖不识趣呢？他说，“历代知识阶层可以制约皇权，制约皇帝，叫做‘从道不从君’。”说白了，就是想要一个传统知识分子向往的理想社会：锦衣玉食地供养着，还要有制约权力的地位。

纵观世界各国的变革历史，政治权力都有一个由皇权——贵族权资产权——官僚精英权的圈子不断扩大的过程，而在这个过程中占人口绝大多数的普通劳动者从来都是排除在权力外的。袁伟时说的“主权在民”，要求分享权力，这个民是不包括劳动人民的。所谓的民主自由只能到此为止，不能再扩大了，所以才有“历史终结论”。意思是资本主义社会已经是人类社会的最高形态了，民主自由只配他们享有，不能更多了，再扩大他们去统治谁呢？所以他才会去丑化辛亥革命，敌视一切历史上的农民起义，鼓吹少数人操控的“温和改革”。在袁伟时看来，民主自由少了不行，多了更不行。

然而历史就是历史，历史不可能按少数人的意愿重新来过。“世界潮流浩浩荡荡，顺之者昌，逆之者亡。”孙中山是一定要革命的，只不过接受了以往的教训。他从 1920 年开始与苏俄人士接触，希望从俄国革命中找到经验。此前列宁对中国的辛亥革命高度关注，曾发表多篇文章，评价了这场革命的世界意义。孙中山在反对北洋政府的斗争中，由于不反帝而幻想争取帝国主义的援助，不依靠国内人民大众而依靠少数帮会、地方军阀，所以屡起屡折，到 1922 年，由他拨巨资支持的军阀陈炯明投靠直系军阀而倒戈。孙中山哀叹“国民党正在堕落中死亡，要救活它需要新血液”。因此，孙中山也欢迎苏俄的援助。1923 年1 月18 日，苏联驻华代表越飞来到上海访晤孙中山后，他们联名发表了《孙文越飞宣言》。同年 8 月 2 日，俄共任命鲍罗廷以共产国际代表名义担任中华民国政府的政治顾问，扶助孙中山改组国民党，在广州建立国民党大元帅府。毛泽东评论孙中山“曾经无数次地向资本主义国家呼吁过援助，结果一切落空，反而遭到了无情的打击。在孙中山的一生中，只得过一次国际的援助，这就是苏联的援助。”

1923 年 12 月 29 日，孙中山落实接受列宁和共产国际的协助重建大元帅府，以苏共为模式重组中国国民党。1924 年 1 月在中国国民党第一次全国代表大会上宣布实行“联俄联共、辅助农工”政策，重新解释三民主义。与共产党合作是一个转折点，自此中国才有了机会摆脱军阀混战的局面。对于选择联俄，孙中山曾有过一段比喻，他说“我们在河中被急流冲走……这时候漂来苏俄这根稻草”。1924 年1 月，国民党一大在广东高等师范学校（后来的中山大学）召开，标志着第一次国共合作的正式形成。这天下午第七项议程时，一个年轻人就“组织国民党政府之必要”作了发言，这个发言引起了孙中山和许多国民党人的关注。孙中山亲自批准这个共产党人为中国国民党章程审查委员。孙中山对这个年轻人赞赏有加，但也没料想正是他接续了辛亥的未竟事业，使反帝反封建的梦想终成现实。这个年轻人，就是毛泽东。